Consumer lending depends on one question that has to be answered before any money moves: how likely is this applicant not to repay? Peer-to-peer (P2P) platforms such as Lending Club publish enough data about their past loans to let students study that question with real, large-scale records rather than toy examples. This report documents a complete, reproducible statistical-learning workflow on the public Lending Club data: it frames the prediction task, builds a leakage-controlled and time-respecting experimental protocol, compares a transparent baseline with several statistical and machine-learning models, and reports honestly where the models help and where they do not.

A guiding rule of the project is that every number and every claim about the data in this report must be produced by code that was run on the data, and that results which disappoint are reported as such. The experiments live in five executed Jupyter notebooks that accompany the report; the tables and figures here are taken from their saved outputs.

We would like to thank our supervisor \todo{supervisor's name} for guidance throughout the course, and the lecturers of the School of Technology for the statistical-learning foundations on which this work rests. Because of limited time and experience, the report may still contain shortcomings; we welcome comments from lecturers and readers so that it can be improved.
